\section{Introduction}\label{sec:intro}

Parasite persistence depends on the hosts needed to complete the life cycle. Complex cycles require two or more host species in sequence and include major agents of human, livestock, and wildlife disease. Examples include mosquito-borne protozoans and filariae such as \textit{Plasmodium falciparum} (human malaria parasite) and \textit{Onchocerca volvulus} \citep{smith2007,basanez2009}, snail-borne trematodes such as \textit{Schistosoma mansoni} and \textit{Schistosoma haematobium} \citep{colley2014,anderson2015}, and tick-borne spirochetes such as \textit{Borrelia burgdorferi} sensu lato \citep{dunn2013,tufts2021}. Trophically transmitted cestodes such as \textit{Echinococcus multilocularis} follow the same obligate pattern across carnivore and rodent hosts \citep{ishikawa2003,atkinson2013}. In all of these systems, at least one required host or vector undergoes wide seasonal or interannual swings in abundance. For instance, rainfall drives mosquito cohorts \citep{ndiath2012,ototo2015,whittaker2022}, hydrological regimes alter freshwater snail populations \citep{rabone2019,anderson2021}, and mast seeding or multi-year cycles govern forest and grassland rodents \citep{elias2004,lambin2006,jacob2011}. Predicting persistence requires knowing whether passage through several hosts buffers those fluctuations or compounds the risk of failure.

Because of this dependence, \citet{dobson2020} asked what governs the dynamics of species with complex life cycles and suggested that using several hosts ``may provide'' ``an intrinsic form of stability''. That chapter advances two distinct arguments for this view. First, a single loop of $n$ stages can be written as a next-generation matrix (the matrix whose entries give the expected infections in the next stage produced per infection in the current stage). Its dominant eigenvalue is the $n$th root of the product of the stage-to-stage contributions, giving a square root for malaria ($n = 2$), a cube root for a trematode ($n = 3$), and a fourth root for a cestode ($n = 4$), or the $n$th root of a route sum when alternative hosts share a stage \citep{dobson2020}. \citet{dobson2020} noted that taking a higher root produces a number ``increasingly close to unity'' and argued that complex-cycle parasites rarely achieve ``explosive birth rates'' ``and thus'' show ``rather stable dynamics''. He added that he could ``only partly understand'' this effect and left a general multi-host solution as a ``conjecture'' and an ``unsolved problem'' \citep{dobson2020}. Second, \citet{dobson2020} suggested that a parasite with several hosts can wait in one long-lived host as a ``parking lot'', such as humans holding \textit{Plasmodium} while mosquitoes are ``transiently reduced''. That storage could allow a parasite to persist where single-host species die out during host ``boom--bust cycles'' \citep{dobson2020}.

Assessing these claims requires separating three meanings of stability, not defined in \citet{dobson2020}, whose threshold calculations use fixed host abundances, although its trophic model includes host and prey differential equations. The chapter does not derive the storage claim from an explicit stochastic-host model. First, invasion and reproductive persistence concern a rare parasite that does not affect its fluctuating environment. We measure this by the sign of the long-run logarithmic growth rate $r$ and the loss in $\ln R_0$ ($R_0$ is the expected descendants completing one full cycle from one individual at baseline). This first reading addresses the storage or ``parking lot'' argument, although a linear invasion exponent is neither an endemic equilibrium nor an extinction probability. Second, stability can mean persistence of an integer-valued population over a finite horizon (or avoidance of a specified decline in a continuous approximation). Finite-population survival depends on offspring clumping, shared parasite death when an infected host dies, starting numbers, and the observation clock in addition to the linear threshold. Third, stability can mean local asymptotic stability of a positive equilibrium (all linearized eigenvalues having negative real parts), either when density dependence regulates loop transfers or when the parasite depresses dynamic host populations \citep{anderson1978,dobson1992}. This reading addresses the $n$th-root argument. Epidemic speed and predator switching in food webs are addressed only in part. Additionally, life-cycle stages and host species are not interchangeable because a free-living egg or larval stage adds a compartment without adding a host species.

For a rare structured population in a fluctuating environment, the mathematical foundation lies in stochastic demography and random matrix products \citep{furstenberg1960,cohen1979b,tuljapurkar1990}. \citet{ruelle1979} proved real-analyticity of the dominant Lyapunov exponent (long-run logarithmic growth rate) under bounded stationary ergodic perturbations of a primitive matrix and gave its second derivative as a lag sum. \citet{lepage1989} and \citet{hennion1991} established differentiability for independent unbounded matrices under invertibility conditions. In population ecology, \citet{tuljapurkar1982a} and \citet{tuljapurkar2006} developed the small-noise growth expansion for serially correlated environments, \citet{tuljapurkar1982b} summed the geometric-memory term over projection-matrix eigenvalues without applying it to multi-host cycles, and \citet{tuljapurkar1985} treated cyclical forcing. Subsequent demographic studies showed that longer generation times buffer variance penalties and that temporal autocorrelation can take either sign \citep{lande1988,orzack1989,morris2008,tuljapurkar2009,engen2013,monmarche2025}. For a two-stage cycle with independent gamma recruitment, \citet{roerdink1988,roerdink1989a} derived the exact stationary ratio law and confluent hypergeometric growth rate in closed form, while \citet{jonsson2021} examined elasticities in two-stage cycles. On the scale of reproduction numbers, \citet{cohen1979b} and \citet{tuljapurkar1990} showed that independent-in-time matrices cannot outgrow their mean matrix. \citet{bacaer2007,bacaer2009} analyzed $R_0$ in periodic models, and \citet{bacaer2013} characterized $R_0$ in finite-state Markov environments as the birth divisor giving zero long-run growth. Parallel genealogical work identified transition elasticities as stationary traversal frequencies along a reproductive-value reference lineage \citep{vangroenendael1994,georgii2003,baake2007,bienvenu2015,bienvenu2017}, proved a reference-chain Jensen bound in random environments \citep{dombry2011} and supplied random-environment path and variational representations \citep{kobayashi2015,kobayashi2017}, and linked fixed-environment growth to cumulants of lineage fitness \citep{nozoe2017,levien2020a,genthon2021,yamauchi2022}.

In addition to demographic perturbation theory, studies of regulated feedback loops, seasonal vectors, and parasite extinction risk bear on both parts of Dobson's argument. In deterministic systems, stability criteria for cyclic feedback networks with equal or unequal decay rates are known in mathematical biology and control theory \citep{blythe1984,lloyd2001,sontag2006,arcak2006,arcak2008,nguyen2009,strelkowa2010}. Classic macroparasite models also show that external larval delays can destabilize host--parasite equilibria, while arrested within-host stages can either stabilize or destabilize them depending on host regulation \citep{anderson1978,may1978,dobson1992,rosa2002,gaba2008,marion2000}. In variable parasite environments, earlier dynamic models analyzed two-host loops under random two-state switching \citep{watts2021}, seasonal malaria persistence via long human infections or liver-stage relapse \citep{smith2007,white2016}, and variance reduction in covert or temperate stages \citep{sorrell2009,maslov2015}. Related work examined statistical averaging across parallel components \citep{doak1998} and thermal shelter inside homeothermic intermediate hosts \citep{molnar2013,kutz2014,aleuy2020}. In contrast, comparative and food-web studies largely associate complex life cycles with heightened vulnerability rather than built-in stability. In presence--absence food-web models, each required host stage adds a point of secondary extinction when hosts disappear \citep{lafferty2009,lafferty2012}. Similarly, field surveys in freshwater lakes tie parasite variability to current-host variability \citep{lagrue2015}, and a century-long museum reconstruction in Puget Sound fishes found steep declines restricted to parasites requiring three or more obligate hosts \citep{wood2023}.

However, the available theory does not by itself decide whether adding a host protects a parasite under a specified comparison. The reference baseline, birth convention, residence-time distribution, and environmental dependence must be fixed. The objective of this study is to test the birth-counting and storage arguments using stage-structured models. We combine residence-time formulas, bounds at every noise level, an exact two-stage benchmark, and local stability criteria. The analysis distinguishes protection through residence times and environmental sampling from the effects of host feedback.

When reproduction is measured per unit time rather than per step of the next-generation matrix, the $n$th-root effect disappears in both invasion thresholds and deterministic regulation (Section~\ref{sec:model}, Proposition~\ref{prop:ring-reg}). Across fluctuating environments, obligate hosts in series do not protect a rare parasite through their number, whereas a long stay in a sheltered stage lowers the second-order fluctuation loss under stated conditions (Section~\ref{sec:ring}; Section~\ref{sec:bounds} bounds the loss at any noise level). In the exactly solved model with independent gamma hosts, a one-step exposed stage, and a geometric steady stay, lengthening that stay strictly reduces the exact threshold loss at every fixed positive variance (Section~\ref{sec:exact}, Theorem~\ref{thm:exact-mono}). Alternative hosts used in parallel at one step behave differently. With equal shares, matched log variance and memory, and fixed mean total transfer, independent hosts divide that step's second-order loss by their number. Partial synchrony weakens this, perfect synchrony removes it, and a much noisier newcomer can raise the loss (Proposition~\ref{prop:ring-alt}). Where stage-based buffering stops or reverses is set by host memory, stage-duration shape, waiting mortality during host scarcity, linked or cyclic host fluctuations, and temporary host absence (Section~\ref{sec:mech}). Once an established parasite depresses its hosts, stability becomes a question of damping around a positive endemic equilibrium (Section~\ref{sec:dyn}), while open empirical archives cannot yet isolate exogenous host variability at fixed mean host abundance and motivate a controlled host-supply experiment (Section~\ref{sec:num}).

Section~\ref{sec:example} works through these ideas with one, two, and three hosts before the general theory begins. Readers mainly interested in the biology can read Sections~\ref{sec:model-ring}--\ref{sec:model-r0} (the proofs can be skipped), Section~\ref{sec:example}, the ``In brief'' openings of Sections~\ref{sec:ring}--\ref{sec:dyn}, and then Sections~\ref{sec:num} and~\ref{sec:disc}.
